\documentclass[12pt]{article}
\usepackage[utf8]{inputenc}
\usepackage[T1]{fontenc}
\usepackage{hyperref}
\title{Tima Language Resources for K-12 Curriculum Development}
\author{Compiled by Claude Sonnet 5 (claude-sonnet-5) for Andr\'as Kornai}
\date{}
\begin{document}
\maketitle

\section{Introduction}
Tima (ISO 639-3 tms; Katla-Tima family; Glottolog places it apart from the Niger-Congo label used in the brief) is spoken in the Nuba Mountains of Sudan, between Katla and Lagawa. Wikipedia lists 3,300 native speakers (2000) and 5,000 ethnic Tima (2007). A 2024 IPA article puts speakers at roughly 7,000. UNESCO classifies Tima as severely endangered. Speakers call it du-murik. The language is unwritten.

\section{Primary Educational Resources}
No primers or school materials were found.

\section{Linguistic and Reference Materials}
\begin{itemize}
\item Wikipedia, ``Tima language''. \url{https://en.wikipedia.org/wiki/Tima_language}
\item ``Tima'', Journal of the International Phonetic Association (Cambridge, 2024), an IPA illustration of the sound system (21 consonants, two tones, twelve vowels with ATR contrasts). \url{https://www.cambridge.org/core/journals/journal-of-the-international-phonetic-association/article/tima/61F761BD22931D58598679FF1F8CD0B1}
\item DoBeS poster on Tima (MPI Nijmegen). The search did not show its contents. \url{https://dobes.mpi.nl/posters/print/africa-tima-print.pdf}
\end{itemize}

\section{Religious and Community Resources}
\begin{itemize}
\item Joshua Project reports that no scripture exists in Tima, though a Bible translation has been started. \url{https://joshuaproject.net/people_groups/15490}
\end{itemize}

\section{Digital and Audio Resources}
None confirmed beyond the DoBeS documentation project poster.

\section{Curriculum Development Context}
Tima is unwritten and endangered, so any curriculum would start with documentation and an agreed orthography. The IPA illustration is the best starting point for the sound system. The DoBeS project may hold recordings, which the search did not confirm.

\end{document}